\documentclass[11pt,a4paper]{article}
\usepackage[utf8]{inputenc}
\usepackage[vietnamese]{babel}
\usepackage{amsmath,amssymb}
\usepackage{geometry}
\geometry{a4paper, total={180mm,260mm}, left=15mm, right=15mm, top=18mm, bottom=18mm}
\usepackage{tikz}
\usetikzlibrary{shapes.geometric, arrows.meta, positioning, calc}
\usepackage{tcolorbox}
\tcbuselibrary{skins,breakable}
\usepackage{xcolor}
\usepackage{multicol}
\usepackage{array}
\usepackage{fontawesome}
\usepackage{fancyhdr}

\pagestyle{fancy}
\fancyhf{}
\renewcommand{\headrulewidth}{0.5pt}
\renewcommand{\footrulewidth}{0.5pt}
\fancyhead[L]{\small\bfseries\color{teal}\faBook\ SỔ TAY KIẾN THỨC TOÁN 3}
\fancyhead[R]{\small\bfseries\color{teal}Lớp Toán Thầy Lĩnh -- loptoanthaylinh.vn}
\fancyfoot[C]{\small\color{gray}-- Sổ tay kiến thức: Bài 32. Mi-li-lít (ml) --}

\definecolor{primary}{HTML}{005F73}
\definecolor{secondary}{HTML}{0A9396}
\definecolor{accent}{HTML}{EE9B00}
\definecolor{danger}{HTML}{AE2012}
\definecolor{boxbg}{HTML}{F4F6F8}

\begin{document}

% --- BANNER TIÊU ĐỀ ---
\begin{tcolorbox}[
  enhanced,
  colback=white,
  colframe=primary,
  arc=3mm,
  boxrule=1.2pt,
  left=4mm, right=4mm, top=2.5mm, bottom=2.5mm,
  drop shadow=black!10!white
]
  \centering
  {\Large\bfseries\color{primary}\faBookmark\ BÀI 32: MI-LI-LÍT (ml)}\\[1.5mm]
  {\small\itshape\color{black!75}Đơn vị đo dung tích nhỏ, quan hệ 1 l = 1000 ml và kĩ năng đọc bình chia vạch trên 1 trang A4}
\end{tcolorbox}

\vspace{-2mm}


\noindent
\begin{minipage}[t]{0.49\textwidth}\vspace{0pt}
\begin{tcolorbox}[
  enhanced,
  colback=blue!4!white,
  colframe=blue!75!black,
  arc=2mm, boxrule=0.8pt,
  left=2mm, right=2mm, top=1.8mm, bottom=1.8mm,
  title={\footnotesize\bfseries\color{white}\faKey\ 1. Đơn vị Mi-li-lít (ml)}
]
{\footnotesize
\begin{itemize}\setlength{\itemsep}{1.5pt}\setlength{\leftskip}{-2.5mm}
  \item Mi-li-lít là đơn vị đo dung tích chất lỏng, viết tắt là \textbf{ml}.
  \item Dùng đo lượng chất lỏng nhỏ: hộp sữa chua, chai thuốc si-rô, cốc nước...
  \item \textbf{Mối quan hệ cốt lõi với Lít ($l$):}
  $$\mathbf{1\ l = 1000\text{ ml}} \iff \mathbf{1000\text{ ml} = 1\ l}$$
  \item \textbf{Phân số dung tích hay gặp:}
  \begin{itemize}\setlength{\itemsep}{0.5pt}\setlength{\leftskip}{-2.5mm}
    \item Nửa lít: $\dfrac{1}{2}\ l = 1000 : 2 = \mathbf{500\text{ ml}}$.
    \item Một phần tư lít: $\dfrac{1}{4}\ l = \mathbf{250\text{ ml}}$.
  \end{itemize}
\end{itemize}
}
\end{tcolorbox}
\end{minipage}
\hfill
\begin{minipage}[t]{0.49\textwidth}\vspace{0pt}
\begin{tcolorbox}[
  enhanced,
  colback=teal!4!white,
  colframe=teal!75!black,
  arc=2mm, boxrule=0.8pt,
  left=2mm, right=2mm, top=1.8mm, bottom=1.8mm,
  title={\footnotesize\bfseries\color{white}\faSliders\ 2. Kĩ năng đọc bình chia vạch}
]
{\footnotesize
\begin{itemize}\setlength{\itemsep}{1.5pt}\setlength{\leftskip}{-2.5mm}
  \item \textbf{Xác định giá trị 1 vạch chia:}\\
  Lấy khoảng cách giữa 2 số ghi chia cho số khoảng nhỏ.
  \item \textbf{Quy tắc đặt mắt đọc số đo:}\\
  Đặt bình trên mặt phẳng nằm ngang, giữ tầm mắt nhìn \textbf{ngang bằng} với mặt thoáng chất lỏng.
  \item Đọc số đo ở vạch ngang trùng với mặt chất lỏng và viết kèm đơn vị $\text{ml}$.
\end{itemize}
}
\end{tcolorbox}
\end{minipage}

\vspace{-2.5mm}
% --- KHỐI 3: SƠ ĐỒ TƯ DUY TIKZ ---
\noindent
\begin{tcolorbox}[
  enhanced,
  colback=orange!4!white,
  colframe=orange!80!black,
  arc=2mm, boxrule=0.8pt,
  left=2mm, right=2mm, top=1.5mm, bottom=1.5mm,
  title={\footnotesize\bfseries\color{white}\faSitemap\ SƠ ĐỒ TƯ DUY: HỆ DUNG TÍCH VÀ QUY TẮC ĐỔI ĐƠN VỊ}
]
\centering
\begin{tikzpicture}[
  box/.style={draw, rounded corners=2pt, align=center, font=\scriptsize\bfseries, inner sep=2.5pt, minimum height=0.75cm},
  arr/.style={-{Stealth[length=2mm]}, thick}
]
  \node[box, fill=purple!15, draw=purple!80!black, text width=2.6cm] (m1) at (0, 0) {ĐƠN VỊ MI-LI-LÍT\\Viết tắt: $\mathbf{ml}$\\Dung tích nhỏ};
  \node[box, fill=teal!15, draw=teal!80!black, text width=2.7cm] (m2) at (3.5, 0) {QUY ĐỔI CHUẨN\\$\mathbf{1\ l = 1000\text{ ml}}$\\$500\text{ ml} = \frac{1}{2}\ l$};
  \node[box, fill=blue!12, draw=blue!80!black, text width=2.6cm] (m3) at (7.0, 0) {ĐỌC CA ĐONG\\Mắt nhìn ngang vạch\\Đếm số khoảng chia};
  \node[box, fill=orange!15, draw=orange!80!black, text width=2.6cm] (m4) at (10.5, 0) {GIẢI TOÁN\\Cộng, trừ, nhân, chia\\kèm tên đơn vị $\text{ml}$};

  \draw[arr, purple!80!black] (m1) -- (m2);
  \draw[arr, teal!80!black] (m2) -- (m3);
  \draw[arr, blue!80!black] (m3) -- (m4);
\end{tikzpicture}
\end{tcolorbox}

\vspace{-2.5mm}
% --- KHỐI 4: VÍ DỤ & MẸO TRÁNH BẪY ---
\noindent
\begin{minipage}[t]{0.49\textwidth}\vspace{0pt}
\begin{tcolorbox}[
  enhanced,
  colback=green!4!white,
  colframe=green!60!black,
  arc=2mm, boxrule=0.8pt,
  left=2mm, right=2mm, top=1.8mm, bottom=1.8mm,
  title={\footnotesize\bfseries\color{white}\faCheckCircle\ Ví dụ thực hành chuẩn mực}
]
{\footnotesize
\textbf{1. Đổi đơn vị dung tích:}\\
$3\ l = \mathbf{3000\text{ ml}}$.\\
$1\ l\ 400\text{ ml} = 1000\text{ ml} + 400\text{ ml} = \mathbf{1400\text{ ml}}$.\\
$2500\text{ ml} = \mathbf{2\ l\ 500\text{ ml}}$.\\
\textbf{2. Bài toán dung tích thực tế:}\\
Một chai nước khoáng có $500\text{ ml}$ nước. Bạn Nam uống hết $200\text{ ml}$. Chai nước còn lại:\\
$500 - 200 = \mathbf{300\text{ ml}}$.
}
\end{tcolorbox}
\end{minipage}
\hfill
\begin{minipage}[t]{0.49\textwidth}\vspace{0pt}
\begin{tcolorbox}[
  enhanced,
  colback=red!4!white,
  colframe=red!70!black,
  arc=2mm, boxrule=0.8pt,
  left=2mm, right=2mm, top=1.8mm, bottom=1.8mm,
  title={\footnotesize\bfseries\color{white}\faExclamationTriangle\ Mẹo tránh bẫy trắc nghiệm}
]
{\footnotesize
\begin{itemize}\setlength{\itemsep}{1.5pt}\setlength{\leftskip}{-2.5mm}
  \item \textbf{Bẫy viết hoa chữ cái lít ($l$):}\\
  Kí hiệu của lít là chữ thường $l$ hoặc nét cong đặc thù, mi-li-lít là $\text{ml}$. Không nhầm sang đơn vị độ dài $\text{mm}$.
  \item \textbf{Bẫy không đổi trước khi thực hiện phép tính:}\\
  $2\ l - 350\text{ ml} \ne 200 - 350$!\\
  Phải đổi $2\ l = 2000\text{ ml} \implies 2000 - 350 = \mathbf{1650\text{ ml}}$.
\end{itemize}
}
\end{tcolorbox}
\end{minipage}


\end{document}
